
i)
Comme \(\underline{O}_Y\) peut se résoudre localement par un complexe de Koszul, l'algèbre graduée \(\bigoplus_{i\ge0}\underline{\Tor}_i^{\underline{O}_X}(\underline{O}_{X'},\underline{O}_Y)\) est alternée (2.2). Il existe donc un homomorphisme d'algèbres et un seul:
\begin{equation}\tag{3.2.1}
\bigwedge \underline{\Tor}_1^{\underline{O}_X}(\underline{O}_{X'},\underline{O}_Y)
\longrightarrow
\bigoplus_{i\ge0}\underline{\Tor}_i^{\underline{O}_X}(\underline{O}_{X'},\underline{O}_Y),
\end{equation}
prolongeant l'injection de \(\underline{\Tor}_1^{\underline{O}_X}(\underline{O}_{X'},\underline{O}_Y)\) dans
\(\bigoplus_{i\ge0}\underline{\Tor}_i^{\underline{O}_X}(\underline{O}_{X'},\underline{O}_Y)\). Il suffit de vérifier localement que c'est un isomorphisme.

On peut donc se placer sur un ouvert affine \(U=\Spec(A)\) de \(X\), \(\underline{J}\) étant défini par un idéal \(J\) de \(A\), et se donner un homomorphisme régulier \(u:E\to A\), définissant \(J\), avec \(E=(A)^n\). Soient \(f_1,\ldots,f_d\) les images d'une base de \(E\).

Dans \(\bigoplus_{n\ge0}J^n\), nous noterons, pour \(m\le p\), \(x^{(m)}\) l'élément \(x\) de \(J^p\) considéré comme étant de degré \(m\). On peut se restreindre à l'ouvert \(V\) de \(X'\) où l'élément \(f_1^{(1)}\) par exemple est inversible.

Les \(\underline{\Tor}_i^{\underline{O}_X}(\underline{O}_{X'},\underline{O}_Y)\) sont les faisceaux d'homologie du complexe \(K_\bullet(u)\otimes_{\underline{O}_X}\underline{O}_{X'}\). Pour étudier ce complexe sur \(V\), on fait le changement de base suivant dans \(E'=E\otimes_A\Gamma(V,\underline{O}_{X'})\): si \(e_1,\ldots,e_d\) est la base donnée de \(E\), on pose:
\[
e'_1=e_1;\qquad \text{pour }i\ne1,\quad
e'_i=e_i-(f_i^{(1)}/f_1^{(1)})e_1.
\]
Posons \(A'=\Gamma(V,\underline{O}_{X'})\); la différentielle du complexe \(K_\bullet(u)\otimes_A A'\) est donnée par
\[
d(e'_1)=f_1^{(0)};\qquad \text{pour }i\ne1,\quad d(e'_i)=0.
\]
En prenant pour base de \(\bigwedge^k(E')\) les éléments
\(e'_{i_1}\wedge\cdots\wedge e'_{i_k}\) avec \(i_1<..<i_k\) on trouve donc:
\[
\begin{array}{ll}
d(a.e'_{i_1}\wedge\cdots\wedge e'_{i_k})=0
&\text{si }i_1\ne1;\\[3pt]
d(a.e'_{i_1}\wedge\cdots\wedge e'_{i_k})=
af_1^{(0)}.e'_{i_2}\wedge\cdots\wedge e'_{i_k}
&\text{si }i_1=1.
\end{array}
\]
